\section{Scope, Method, and Qualifications}

\subsection*{Subject and Intended Reader}

This study is Part IV of a four-part Catholic series on the Four Last Things (Death, Judgment, Hell, and Heaven). Written for serious Catholic believers, it traces the Church's doctrine on heaven from Sacred Scripture and conciliar definitions through the Fathers, Scholastic Doctors, liturgy, and saints into examination of conscience and received vocal prayer. Questions belonging primarily to physical death and preparation for dying (Part I), Particular and General Judgment and Purgatory (Part II), or eternal damnation (Part III) are touched only where they directly illuminate the kingdom prepared for the just.

\subsection*{Scripture, Liturgical Texts, and Received Prayers}

English quotations of Sacred Scripture follow the Challoner revision of the Douay-Rheims Bible, paired where helpful with the Clementine Vulgate Latin and Vulgate verse numbering throughout (notably Psalms 16, 33, 35, 41, 72, 83, and 102).

Latin liturgical texts follow the 1962 \textit{Missale Romanum} for the Feast of All Saints (collated against the 1922 Tours Mame \textit{editio quarta iuxta typicam}) and the 1604 Vatican typical \textit{Missale Romanum} for the Roman Canon, with English renderings drawn from Eugene Cummiskey's 1861 bilingual Roman Missal (preserving verbatim its plural rendering ``by the hands of thy holy angels'' alongside an endnote recording the singular Latin \textit{sancti Angeli tui}, and noting Pope St.\ John XXIII's November 1962 insertion of St.\ Joseph into the \textit{Communicantes}) and collated with Rev.\ F.\ X.\ Lasance's revised 1945 \textit{New Roman Missal}. Venerable Bede's Matins homily for the Second Day within the Octave of All Saints (2 November) follows the Marquess of Bute's 1879 \textit{Roman Breviary}, explicitly situated in the classical pre-1955 calendar prior to the suppression of the Octave by the Sacred Congregation of Rites' decree \textit{Cum nostra hac aetate} (23 March 1955).

Every text set for vocal prayer in Section 10 reproduces a received public-domain edition verbatim: the Latin and English \textit{Te Deum laudamus} and English All Saints Collect from the 1889 \textit{Manual of Prayers for the Use of the Catholic Laity} (Third Plenary Council of Baltimore), the Latin All Saints Collect from the 1962 \textit{Missale Romanum}, and St.\ Alphonsus Maria de Liguori's prayer on Heaven from the 1843 Dublin Duffy \textit{Preparation for Death} (resolving only obvious nineteenth-century OCR character and spacing glitches without altering wording). Brief translations of Latin phrases inside analytical prose serve solely as explanatory glosses.

\subsection*{Source Hierarchy, Editions, and Limits}

Appendix A distinguishes five levels of authority: Sacred Scripture and solemn definitions (\textit{de fide}); ordinary magisterial teaching; patristic and scholastic consensus (\textit{sententia communis}); open theological questions among approved Catholic Doctors (such as whether bodily reunion increases essential beatitude in intensity or only in extent, how risen eyes perceive God's glory, and the physical mode of cosmic renewal); and approved private testimony (St.\ Teresa of Avila's \textit{Life}, chs.\ 38 and 40).

Patristic, scholastic, and spiritual texts are quoted directly from identified public-domain editions: \textit{Ante-Nicene Fathers}, \textit{Nicene and Post-Nicene Fathers}, Arthur West Haddan's 1887 translation of St.\ Augustine's \textit{De Trinitate} as revised by Kevin Knight, Marcus Dods's 1871 \textit{City of God} (collated with CSEL 40 Latin), the Latin PL 76 text of St.\ Gregory the Great's \textit{Homiliae in Evangelia}, Edward Marshall's 1878 Venerable Bede \textit{Explanation of the Apocalypse}, Marianne Caroline and Coventry Patmore's 1884 St.\ Bernard \textit{De diligendo Deo}, the English Dominican \textit{Summa Theologiae} (noting a minor Gutenberg slip of ``1 John 2:2'' for 1 John 3:2 in \textit{ST} I, q.\ 12, a.\ 5) and posthumously compiled \textit{Supplement} (drawn from St.\ Thomas's \textit{Scriptum super Sententiis} IV, 1922 English Dominican edition covering qq.\ 87--99), J.\ Donovan's 1829 \textit{Roman Catechism} (bracketing ``[acquire]'' where the 1829 printing reads ``require'' for Latin \textit{acquiret} in Creed Art.\ XI), the Lewis--Zimmerman 1904 \textit{Life of St.\ Teresa}, the Pustet Roman Catholic \textit{Introduction to a Devout Life}, and the 1843 Duffy \textit{Preparation for Death}. Among conciliar texts, Fourth Lateran Council Constitution 1 (\textit{Firmiter credimus}, DS 800--801) and First Vatican Council Dogmatic Constitution \textit{Dei Filius} (ch.\ 1, DS 3001) are read directly in their historical Latin texts. Conciliar and papal acts held as cataloged bibliographic records---Council of Vienne \textit{Ad nostrum} (DS 895), Pope Benedict XII's \textit{Benedictus Deus} (DS 1000--1001), Council of Florence \textit{Laetentur caeli} (DS 1305), Council of Trent Session VI Canon 32 (DS 1582), Pope Pius XII's \textit{Munificentissimus Deus} (\S 44, DS 3903) and \textit{Ad Caeli Reginam}, and the \textit{Catechism of the Catholic Church} (\S\S 1023--1029, 1042, 1048)---are cited strictly by canonical DS or section locus in doctrinal paraphrase without verbatim quotation or any claim of local text inspection.

\subsection*{Rights and Non-Authority Notice}

Offered as an independent aid for reading, reflection, and prayer, this work is not an official liturgical book, catechism, magisterial pronouncement, or substitute for a confessor or spiritual director. Project-authored exposition carries the repository license stated in the closing colophon; Scripture, official liturgical and magisterial texts, received prayers, and historical translations retain their independent rights status.
